% instructor/w_intensivo_guia_instructor.tex
\documentclass[11pt, a4paper]{article}

\usepackage[T1]{fontenc}
\usepackage[utf8]{inputenc}
\usepackage{tgpagella}
\usepackage[spanish, es-noshorthands]{babel}
\usepackage{amsmath, amssymb, bm}
\usepackage[most]{tcolorbox}
\usepackage[margin=2.2cm]{geometry}
\usepackage{fancyhdr}
\usepackage{booktabs}
\usepackage{tabularx}
\usepackage{enumitem}

% Definición de columna para evitar Underfull \hbox en tabularx
\newcolumntype{L}{>{\raggedright\arraybackslash}X}

\pagestyle{fancy}
\fancyhf{}
\setlength{\headheight}{16pt}
\lhead{\textbf{UCB Sede Tarija} | Ing. Mecatrónica}
\rhead{IMT-221 | Guía Docente, Clave WE1--WE9 y Gates}
\cfoot{\thepage}
\renewcommand{\headrulewidth}{0.4pt}
\definecolor{ucbblue}{RGB}{0, 51, 102}

\begin{document}

\begin{center}
    {\LARGE \textbf{Guía del Instructor, Solucionario Completo (WE1--WE9) y Protocolo de Gates}}
    \par\vspace{0.15cm}
    \textit{IMT-221 — Laboratorio Intensivo Integrado (10 al 17 de noviembre de 2026) — USO EXCLUSIVO DOCENTE}
\end{center}

\begin{tcolorbox}[colback=ucbblue!5, colframe=ucbblue, title=\textbf{Directriz de Gestión Académica y Seguridad}]
    Esta reorganización pedagógica concentra la recolección de evidencias pendientes en un único proyecto acumulativo entre el 10 y el 16 de noviembre, con defensa prevista el martes 17/11/2026 (sujeta a ratificación institucional). \textbf{No convalida automáticamente sesiones ni checkpoints previos (CP1A, CP1B, CP2A):} cada competencia debe demostrarse con evidencia trazable y defensa individual.
\end{tcolorbox}

\section{Secuencia Operativa y Gates Formativos (09/11 al 17/11)}

\begin{center}
\begin{tabularx}{\textwidth}{@{} l l L L l @{}}
\toprule
\textbf{Fecha} & \textbf{Bloque} & \textbf{Rol del Instructor} & \textbf{Evidencia Exigida al Estudiante} & \textbf{Gate / Hito} \\ \midrule
Lun 09/11 & Prep. & Auditar inventario, datasheets y confirmar reglas institucionales & Lista real de recursos y archivos previos S1--S3 & Inven. s/energía \\ \addlinespace[3pt]
Mar 10/11 & Arranque & Aplicar Gate 1 ($H(s)$) y Gate 2 (Seguridad); fijar Ruta A/B/C & Recuperación S3, predicciones DC, BOM firmado & \textbf{Gate 1, 2} (Ruta) \\ \addlinespace[3pt]
Mié 11/11 & Op-Amp & Verificar pinout/rieles; aplicar Gate 3 (límites reales) & Cálculos SR/GBW con datasheet LT1001 y SPICE & \textbf{Gate 3} (\texttt{CP2B}) \\ \addlinespace[3pt]
Jue 12/11 & Filtro & Aplicar Gate 4 (predicción previa); supervisar montaje Ruta A & Sallen-Key KCL, tabla 4 frec., barrido AC y rediseño & \textbf{Gate 4} (\texttt{CP3A}) \\ \addlinespace[3pt]
Vie 13/11 & Control & Aplicar Gate 5 (P vs PI) y \textbf{Freeze del montaje físico} & Comparación P/PI/PID Python, análisis windup/ruido & \textbf{Gate 5} (\texttt{CP4A}) \\ \addlinespace[3pt]
Sáb 14/11 & Integrac. & Aplicar Gate 6 (cambio de parámetro no ensayado) & Microcasos Wien/AB (\texttt{CP3B/4B}), tabla antes/después & \textbf{Gate 6} (\texttt{CP5}) \\ \addlinespace[3pt]
Dom 15/11 & Redacción & Facilitar plantilla IEEE y auditar leyendas de procedencia & Borrador completo IEEE sin datos ficticios & \textbf{Gate 7} (Borrador) \\ \addlinespace[3pt]
Lun 16/11 & Cierre & Mock tribunal, auditoría de carpeta reproducible y preguntas & Informe final IEEE, carpeta y defensa ensayada & Cierre continuo \\ \addlinespace[3pt]
Mar 17/11 & Examen & Tribunal evalúa demostración, defensa y rediseño en vivo & Demostración delimitada, defensa y desafío \textbf{E11} & Examen / Acta \\ \bottomrule
\end{tabularx}
\end{center}

\section{Solucionario Técnico Completo de Ejemplos Desarrollados (WE1 a WE9)}

\subsection*{WE1: Mapeo Lineal de Señal y Acondicionador (Bloque 1 y E1)}
\begin{itemize}[noitemsep]
    \item \textbf{Mapeo sin offset:} Puntos $(0.20, 0.60)\,\text{V}$ y $(0.80, 2.40)\,\text{V}$:
    \[
        a = \frac{2.40 - 0.60}{0.80 - 0.20} = 3\,\text{V/V}, \qquad b = 0.60 - 3(0.20) = 0\,\text{V}
    \]
    Con $R_g = 10\,\text{k}\Omega$ y $R_f = 20\,\text{k}\Omega$, $A_v = 1 + \frac{20}{10} = 3$. Entradas $0.20, 0.50, 0.80\,\text{V}$ producen salidas ideales $0.60, 1.50, 2.40\,\text{V}$.
    \item \textbf{Variaciones de E1:} Si $V_s = 0.90\,\text{V} \implies V_{out,\text{nom}} = 2.70\,\text{V}$. Si $V_s = 1.50\,\text{V} \implies V_{out,\text{nom}} = 4.50\,\text{V}$ (requiere verificar que el swing del op-amp bajo la alimentación real permita alcanzar $4.50\,\text{V}$ sin saturar). Si $R_f$ cambia a $30\,\text{k}\Omega$, $A_v = 1 + \frac{30}{10} = 4\,\text{V/V}$ y para $0.80\,\text{V}$ la salida ideal pasa a $3.20\,\text{V}$.
\end{itemize}

\subsection*{WE2: Red RC de Primer Orden (Bloque 2 y E0)}
Con $R = 10\,\text{k}\Omega$ y $C = 100\,\text{nF}$:
\[
    \tau = RC = 10^4 \times 10^{-7} = 1\,\text{ms}, \qquad f_c = \frac{1}{2\pi(0.001)} \approx 159.155\,\text{Hz}
\]
En $f = f_c$, $|H(j2\pi f_c)| = \frac{1}{\sqrt{2}} \approx 0.7071$ ($-3.0103\,\text{dB}$) y $\phi(f_c) = -45^\circ$. Si $C$ se duplica a $200\,\text{nF}$, $\tau' = 2\,\text{ms}$ y $f_c' \approx 79.577\,\text{Hz}$.

\subsection*{WE3: Calificación de Slew Rate y Ancho de Banda del LT1001 (Bloque 3 y E2)}
Bajo condiciones de ficha ($V_S = \pm 15\,\text{V}, R_L \ge 2\,\text{k}\Omega$), el LT1001 especifica $GBW_{\min} = 0.4\,\text{MHz}$ (típ. $0.8\,\text{MHz}$) y $SR_{\min} = 0.1\,\text{V/}\mu\text{s}$ (típ. $0.25\,\text{V/}\mu\text{s}$). Para una salida sinusoidal con $V_p = 2\,\text{V}$:
\begin{align*}
    f = 1\,\text{kHz}: \quad SR_{\text{req}} &= 2\pi(1000)(2) \approx 12566\,\text{V/s} \approx 0.0126\,\text{V/}\mu\text{s}
\end{align*}
\textit{Conclusión:} Cumple con amplio margen frente al mínimo $0.1\,\text{V/}\mu\text{s}$.
\begin{align*}
    f = 10\,\text{kHz}: \quad SR_{\text{req}} &= 2\pi(10000)(2) \approx 125664\,\text{V/s} \approx 0.1257\,\text{V/}\mu\text{s} 
\end{align*}
\textit{Conclusión:} Viola el mínimo garantizado de $0.1\,\text{V/}\mu\text{s}$. \\
\textbf{Criterio docente:} Aunque $10\,\text{kHz}$ es muy inferior al ancho de banda de pequeña señal estimado ($f_{BW,\min} \approx 400\,\text{kHz}/3 \approx 133.3\,\text{kHz}$), el dispositivo no puede garantizar operación libre de distorsión por pendiente a $2\,\text{V}$ pico porque supera el $SR_{\min}$ contractual.

\subsection*{WE4 y WE5: Filtro Sallen-Key, Tabla Resuelta y Rediseños (Bloque 4, E4 y E5)}
Con $R_1 = R_2 = 16\,\text{k}\Omega$, $C_1 = C_2 = 100\,\text{nF}$, $R_g = 10\,\text{k}\Omega$ y $R_f = 5.9\,\text{k}\Omega$:
\[
    K = 1 + \frac{5.9}{10} = 1.59 \quad (+4.03\,\text{dB}), \qquad Q = \frac{1}{3 - 1.59} \approx 0.709, \qquad f_0 = \frac{1}{2\pi(1.6\times 10^{-3})} \approx 99.47\,\text{Hz}
\]

\begin{center}
\footnotesize
\begin{tabularx}{\textwidth}{@{} c c X X X X @{}}
\toprule
$\bm{f}$ \textbf{[Hz]} & $\bm{x = f/f_0}$ & $\bm{|H|/K}$ \textbf{[V/V]} & $\bm{G_{\text{relativa}}}$ \textbf{[dB]} & $\bm{|H|}$ \textbf{[V/V]} & $\bm{G_{\text{absoluta}}}$ \textbf{[dB]} / $\bm{\phi}$ \textbf{[$^\circ$]} \\ \midrule
$10.00$ & $0.101$ & $1.000$ & $+0.00\,\text{dB}$ & $1.590$ & $+4.03\,\text{dB} \quad / \quad -8.15^\circ$ \\
$99.47$ & $1.000$ & $0.709$ & $-2.98\,\text{dB}$ & $1.128$ & $+1.04\,\text{dB} \quad / \quad -90.00^\circ$ \\
$1000$ & $10.05$ & $0.0099$ & $-40.09\,\text{dB}$ & $0.0157$ & $-36.06\,\text{dB} \quad / \quad -171.9^\circ$ \\
$10000$ & $100.5$ & $0.0001$ & $-80.09\,\text{dB}$ & $0.00016$ & $-76.06\,\text{dB} \quad / \quad -179.2^\circ$ \\ \bottomrule
\end{tabularx}
\end{center}

\begin{itemize}[noitemsep]
    \item \textbf{Clave de E5 (Cambio de $R_f$ de $5.9\,\text{k}\Omega$ a $10\,\text{k}\Omega$):} $K' = 1 + \frac{10}{10} = 2.0$ ($+6.02\,\text{dB}$). El factor de calidad sube a $Q' = \frac{1}{3 - 2} = 1.0$. El filtro deja de ser Butterworth y presenta un pico de resonancia en $f_{\text{pico}} \approx 70.3\,\text{Hz}$, donde la ganancia relativa alcanza $+1.25\,\text{dB}$ sobre la banda de paso ($+7.27\,\text{dB}$ absolutos).
    \item \textbf{Clave de WE5 (Duplicar capacitores a $200\,\text{nF}$):} $f_0$ baja a $49.7\,\text{Hz}$ sin alterar $Q = 0.709$. Si se modifica \textbf{un solo} componente, se rompe la simetría y debe usarse la ecuación de $Q$ general.
\end{itemize}

\subsection*{WE6, WE7, WE8 y WE9: Control P/PI, Potencia y Contraste de Modelos}
\begin{itemize}[noitemsep]
    \item \textbf{WE6 (Control P, $K_c=1$):} $T_P(s) = \frac{2}{20s+3}$, ganancia DC $2/3$, error estacionario $e_\infty = 1/3$, $\tau_{CL} = 6.67\,\text{s}$.
    \item \textbf{WE7 (Control PI, $K_c=1, K_i=0.05\,\text{s}^{-1}$):} $C_{\text{PI}}(s) = \frac{20s+1}{20s} \implies L(s) = \frac{1}{10s} \implies T_{\text{PI}}(s) = \frac{1}{10s+1}$. Error estacionario $e_\infty = 0$. Si la constante $\tau$ real difiere de $20\,\text{s}$, la cancelación es inexacta y surge un modo lento.
    \item \textbf{WE8 (Potencia):} $5\,\text{V}$ DC sobre $100\,\Omega \implies I = 50\,\text{mA}, P = 250\,\text{mW}$. Excede la capacidad de salida de un op-amp de señal común.
    \item \textbf{WE9 (Modelos SPICE):} El VCVS ideal ($10^6$) tiene banda infinita; al sustituirlo por el macromodelo oficial se observa la caída por el $GBW$ real y el cambio de fase adicional.
\end{itemize}

\newpage
\section{Banco Reservado para el Desafío Imprevisto de Defensa (E11)}
Durante el examen del martes 17/11 (o en el Gate 6 del sábado 14/11), asigne al estudiante \textbf{una} modificación no resuelta en su guía y exija predicción analítica escrita antes de simular:
\begin{enumerate}[noitemsep]
    \item \textbf{Variante E11-A (Asimetría en Sallen-Key):} Manteniendo $R_2 = 16\,\text{k}\Omega, C_1 = C_2 = 100\,\text{nF}, K = 1.59$, cambiar únicamente $R_1$ a $32\,\text{k}\Omega$.
    \begin{itemize}[noitemsep]
        \item \textit{Solución docente:} $f_0' = \frac{99.47}{\sqrt{2}} \approx 70.3\,\text{Hz}$. Usando la fórmula general:
        \begin{align*}
            Q' &= \frac{\sqrt{R_1 R_2} C}{(R_1 + R_2)C + R_1 C(1 - K)} \\
            Q' &= \frac{\sqrt{32 \times 16}}{32 + 16 + 32(1 - 1.59)} \approx \frac{22.63}{29.12} \approx 0.777
        \end{align*}
        Se reduce $f_0$ y aumenta $Q'$, superando el valor Butterworth.
    \end{itemize}
    \item \textbf{Variante E11-B (Seguidor Sallen-Key Unitario):} Reemplazar el bloque $K=1.59$ por un seguidor unitario ($K = 1$, es decir $R_f = 0, R_g \to \infty$) manteniendo $R_1 = R_2 = 16\,\text{k}\Omega, C_1 = C_2 = 100\,\text{nF}$.
    \begin{itemize}[noitemsep]
        \item \textit{Solución docente:} $f_0$ sigue en $99.47\,\text{Hz}$, pero $Q' = \frac{1}{3 - 1} = 0.50$ (críticamente amortiguado). En $f_0$, la ganancia relativa ya no es $-3.0\,\text{dB}$ sino $20\log_{10}(0.5) = -6.02\,\text{dB}$.
    \end{itemize}
    \item \textbf{Variante E11-C (Perturbación de Ganancia en Planta):} Si la ganancia estática de la planta se reduce a la mitad ($G_p'(s) = \frac{1}{20s+1}$) manteniendo $C_P(s)=1$ y $C_{\text{PI}}(s) = \frac{20s+1}{20s}$.
    \begin{itemize}[noitemsep]
        \item \textit{Solución docente:} En control P, $T_P'(s) = \frac{1}{20s+2} \implies y(\infty) = 0.5, e_\infty = 0.5, \tau_{CL}' = 10\,\text{s}$. En control PI, $L'(s) = \frac{1}{20s} \implies T_{\text{PI}}'(s) = \frac{1}{20s+1}$, manteniendo $e_\infty = 0$ pero duplicando la constante de tiempo a $20\,\text{s}$.
    \end{itemize}
\end{enumerate}

\section{Lista de Verificación Post-Clase y Auditoría Docente (17 Preguntas)}
Al cierre de cada jornada y antes del examen del 17/11, el instructor debe verificar:
\begin{enumerate}[noitemsep, font=\footnotesize]
    \item ¿Qué sesiones se registraron administrativamente cada día y cuáles no?
    \item ¿Qué ruta (A, B o C) fue autorizada y sobre qué inventario físicamente comprobado?
    \item ¿Qué circuitos se construyeron en hardware, cuáles se simularon y cuáles no se hicieron?
    \item ¿Qué modelo exacto de op-amp, alimentación e instrumentos quedaron documentados?
    \item ¿Se confirmó y superó la dificultad histórica con $H(s)$ mediante predicción individual?
    \item ¿Qué evidencias previas (\texttt{CP1A}, \texttt{CP1B}, \texttt{CP2A}) se revisaron y cuál es su estado formal?
    \item ¿Qué nuevos checkpoints tienen pruebas suficientes y cuál es su alcance (\texttt{MEDIDO}/\texttt{SIMULADO}/mixto)?
    \item ¿Los datos térmicos son reales o ilustrativos, y cómo se declara el modelo de planta?
    \item ¿La cadena física (si aplica Ruta A/B) opera dentro de rangos seguros y medibles?
    \item ¿Se registraron los cuatro puntos de frecuencia del filtro y las diferencias teoría/simulación/medición?
    \item ¿Existe al menos un rediseño controlado con predicción previa y evidencia antes/después?
    \item ¿El estudiante distingue rigurosamente lo que construyó de lo que únicamente modeló?
    \item ¿El informe IEEE enlaza todas sus tablas y figuras con archivos fuente reproducibles?
    \item ¿Se confirmó institucionalmente la fecha del examen (17/11), requisitos de habilitación y tribunal?
    \item ¿Se cerró la evaluación continua con trazabilidad documental y sin calificaciones inferidas?
    \item ¿Qué riesgos o criterios siguen abiertos al cierre del 16/11 y cómo se declararán ante el tribunal?
    \item ¿La defensa oral evidenció razonamiento propio ante la modificación no ensayada (\textbf{E11})?
\end{enumerate}

\end{document}
